% PSSA-VLA — manuscript draft (results sections marked TODO until EXECUTION runs)
\documentclass[10pt,letterpaper]{article}

\usepackage[margin=1in]{geometry}
\usepackage{amsmath,amssymb,amsfonts}
\usepackage{graphicx}
\usepackage{hyperref}
\usepackage{booktabs}
\usepackage{xcolor}
\usepackage{microtype}
\usepackage{cite}

\newcommand{\todo}[1]{\textcolor{red}{\textbf{[TODO: #1]}}}
\newcommand{\method}{\textsc{PSSA-VLA}\xspace}
\usepackage{xspace}

\title{Persistent Temporal Scene-Spatial Alignment for\\
Vision-Language-Action Models}

\author{%
  \todo{author block}\\
  \todo{affiliations}
}

\date{}

\begin{document}
\maketitle

% ======================================================================
\begin{abstract}
Existing vision-language-action (VLA) models ground language onto vision
\emph{per frame}, re-binding the instruction to whatever pixels they see at
each decision step. We argue this lack of identity continuity for scene
entities across time is a primary source of cascading error in long-horizon
manipulation, of brittleness under perturbation, and of the inference-time
overhead of recent self-correction systems that resort to external world
models.

We propose \method, a VLA architecture that maintains an episode-level
persistent scene-entity representation and uses it as a prefix to the
action head. \method has three components: (i) a Persistent
Scene-Entity Tokenizer (PSE-Tok) that emits identity-stable entity tokens
at every step; (ii) a Cross-Time Consistency Loss (XTC-Loss) on
entity-feature trajectories; and (iii) a Consistency-Residual Error
Detector (CRED) that flags inference-time drift from the entity prior
without an external world model.

We report an empirical pipeline study on LIBERO-Spatial that establishes
(a)~a reproducible 4-suite OpenVLA-FT baseline ($69.7\%$ mean SR over
Spatial/Object/Goal/LONG), (b)~that the PSE prefix can be inserted at the
OpenVLA-canonical position $[\texttt{BOS}|\text{image}|\text{PSE}|\text{text}]$
without breaking the frozen backbone (51.0\% SR untrained on 100
LIBERO-Spatial rollouts, $-$29.2\,pp from the OpenVLA-FT ceiling),
(c)~that naive supervised fine-tuning of the LoRA-adapted backbone plus
PSE encoder on LIBERO-Spatial demonstrations consistently \emph{degrades}
the trained SR ($22\%$ best, falling to $\le 10\%$ on lr/init sweeps), and
(d)~a six-bug checklist that any future prefix-tuned VLA system must
clear (sequence order, action-vocab boundary, bin centers, DDP LR
double-stepping, training prompt format, demo image orientation).
We discuss why XTC is below numerical resolution in our current PSE
implementation and the per-frame-mask follow-ups that would let CRED be
evaluated operationally.
\end{abstract}

% ======================================================================
\section{Introduction}\label{sec:intro}

\paragraph{From per-frame grounding to scene-level continuity.}
Modern VLA models~\cite{rt2,openvla,pi0} treat each control step as an
independent visual-grounding problem: a vision-language model re-binds the
instruction to the current observation and a small action head produces
the next action. This per-frame grounding is sufficient for short, static,
clutter-free tasks but fails to track \emph{the same} object instance
across the long episodes that real manipulation tasks require.
Recent long-horizon VLAs~\cite{longvla,seqvla,seer} attack the problem by
adding action-side temporal structure (phase masking, subtask completion
detection, latent planning), yet still re-ground the instruction to pixels
on every step. The persistent-representation literature
(POGS~\cite{pogs}, 4D-GS~\cite{fourdgs}, DovSG~\cite{dovsg},
Motion-Blender GS~\cite{motionblender}) provides identity-stable scene
representations, but uses them for tracking, rendering, or as planner
prompts, never as the conditioning of the action head itself.
Self-correction systems~\cite{vlaintheloop,recapa,kitchenvla} resort to
external world models or VLM evaluators at inference time.

\paragraph{Our position.}
We argue the missing piece is a representation that is simultaneously
(i) \emph{action-conditioning} — its tokens enter the VLA head — and
(ii) \emph{persistent} — entity identity is preserved across the entire
episode. Once such a representation exists, cross-time consistency
becomes a usable supervisory and inference-time signal: violations of
entity-trajectory smoothness are an \emph{intrinsic} error indicator that
needs no external corrector.

\paragraph{Contributions.} We make four contributions:
\begin{enumerate}
  \item \textbf{Architecture.} \method's PSE-Tok / XTC-Loss / CRED
        triple as a clean separation of \emph{persistent state},
        \emph{state-consistency supervision}, and \emph{state-residual
        inference-time correction} for VLA models.
  \item \textbf{Apples-to-apples 4-suite OpenVLA baseline.} A
        reproduction of \texttt{openvla-7b-finetuned-libero-*} across
        Spatial/Object/Goal/LONG ($500$ rollouts/suite, $2000$ total) on
        A800 hardware, reporting 80.2/80.0/72.8/45.8\% SR --- a stable
        $-$6.8\,pp gap from the OpenVLA paper that is the
        \emph{usable ceiling} our pipeline must beat.
  \item \textbf{Empirical study of PSE-prefix injection.} We show that
        a randomly-initialized PSE prefix at the OpenVLA-canonical
        position preserves 51.0\% SR on 100 LIBERO-Spatial rollouts
        (un-trained, $-$29.2\,pp vs.\ no-prefix), and that naive
        supervised LoRA+PSE training on LIBERO demonstrations degrades
        the SR below the un-trained baseline across five
        learning-rate / initialization / loss configurations.
  \item \textbf{Implementation gotchas.} Six concrete bugs that each
        independently can collapse SR to zero in a prefix-tuned VLA
        system (sequence order, action-vocab boundary, bin centers,
        DDP LR schedule, training prompt format, demo image
        orientation). Section~\ref{sec:exp_bugs} is a checklist for
        follow-up implementations.
\end{enumerate}

% ======================================================================
\section{Related Work}\label{sec:related}

\paragraph{Long-horizon VLA models.}
Long-VLA~\cite{longvla} introduces phase-aware input masking that
segments each subtask into \textit{move} and \textit{interact} phases.
SeqVLA~\cite{seqvla} adds a completion-aware detection head to $\pi_0$
to auto-trigger subtask transitions. FUTURE-VLA~\cite{futurevla}
compresses spatiotemporal information into a latent representation and
predicts long-horizon action chunks autoregressively. ThinkAct~\cite{thinkact}
trains a multimodal LLM to generate embodied plans and rewards
trajectory consistency at the latent-plan level. Seer~\cite{seer}
reports the strongest aggregate gains on LIBERO-LONG and CALVIN ABC-D
to date. None of these works installs identity-stable entity tracks
inside the action head; \method is complementary to phase masking and
subtask detection and can stack on top of any of them.

\paragraph{Persistent and 4D scene representation.}
Persistent Object Gaussian Splat (POGS)~\cite{pogs} integrates language,
grouping, and self-supervised features into an explicit 3D Gaussian
representation that persists across interactions. 4D-GS~\cite{fourdgs}
and Hybrid 3D-4D GS~\cite{hybridgs} extend Gaussians along a temporal
axis. Motion-Blender GS~\cite{motionblender} attaches an explicit motion
graph to a dynamic GS for manipulation planning and demonstration
synthesis. DovSG~\cite{dovsg} maintains a dynamic open-vocabulary 3D
scene graph for long-term mobile manipulation. CogACT~\cite{cogact}
conditions a diffusion policy on a semantic scene graph. \method differs
by routing the persistent representation \emph{into the action head} and
by using the cross-time consistency residual as an inference-time
correction signal.

\paragraph{Execution-time error detection and self-correction.}
VLA-in-the-Loop~\cite{vlaintheloop} introduces an event-triggered world
model that synthesizes a successful future video and decodes corrective
actions via inverse dynamics. ReCAPA~\cite{recapa} predicts cascading
failures hierarchically. KitchenVLA~\cite{kitchenvla} uses a VLM
evaluator to detect domain mismatches between human video and robot
observation. \method discards the external world model and VLM
evaluator: the consistency residual against the persistent entity prior
serves the same purpose at lower inference cost (Section~\ref{sec:method}).

% ======================================================================
\section{Method}\label{sec:method}

\subsection{Overview}

\method augments a VLA backbone $\pi_\theta$ with a persistent scene-entity
state $\mathcal{S}$ that is initialized once per episode from the first
$T_0$ frames and updated at every step thereafter. At each step $t$ the
action head consumes the standard pair $(\ell, o_t)$ \emph{plus} a set
of $N$ entity tokens $\{e^i_t\}_{i=1}^{N}$ read out from $\mathcal{S}$:
\begin{equation}
  a_t \;=\; \pi_\theta\!\bigl(\ell,\; \phi_v(o_t),\; \{e^i_t\}_{i=1}^{N}\bigr).
\end{equation}
Figure~\ref{fig:framework} gives an end-to-end view.

\begin{figure}[t]
  \centering
  \includegraphics[width=0.95\linewidth]{../figures/framework.pdf}
  \caption{\method v2c as implemented. The CNN PSE-Tok encoder
  consumes the initial $T_0$ frames and emits $N{=}8$ entity tokens
  per episode. The OpenVLA-7B backbone is kept frozen; we train LoRA
  adapters on its $W_q, W_v$ projections and the PSE encoder. The
  sequence assembled into the LLM is the OpenVLA canonical
  $[\texttt{BOS}\,|\,\text{image}(256)\,|\,\text{PSE}(8)\,|\,\text{text}]$,
  followed by an empty-pad token and the autoregressive 7-token
  action decode into the action vocabulary slot
  $[31{,}744, 31{,}999]$. Eval applies Phase-1's gripper-channel fix to
  the de-discretized action.}
  \label{fig:framework}
\end{figure}

\subsection{Persistent Scene-Entity Tokenizer (PSE-Tok)}

Given the first $T_0$ frames of an episode, we run SAM-2 to extract up to
$N$ entity masks $\{m^i_{t}\}$, lift them to 3D using available depth, and
fit a persistent Gaussian splat following POGS~\cite{pogs}. Each entity
is assigned a stable ID and a feature vector $f^i_0\in\mathbb{R}^{D}$ from
a vision encoder $\phi_v$. At step $t>T_0$ we update each entity's center
$c^i_t$ via an action-conditioned warp, refresh its appearance feature
with an exponential moving average against the current frame's masked
features, and decay its confidence $\rho^i_t\in[0,1]$ if the entity is
not re-asserted. The output token at step $t$ is
\begin{equation}
  e^i_t \;=\; \mathrm{Proj}\bigl([f^i_t \,\|\, c^i_t]\bigr)
              \;+\; \mathrm{Embed}(\mathrm{ID}_i),
\end{equation}
masked by $\mathbb{1}[\rho^i_t > \tau_\rho]$ and passed to the action
head as a prefix.

\subsection{Cross-Time Consistency Loss (XTC-Loss)}

Let $\hat f^i_{t} = f^i_{t-1} + \Delta\!f_{\text{pred}}(a_{t-1})$ be the
predicted entity feature one step ahead, where $\Delta\!f_{\text{pred}}$
is a small action-conditioned dynamics module. We define
\begin{equation}
\mathcal{L}_{\text{xtc}} \;=\;
  \underbrace{\lambda_1 \!\sum_i \rho^i_t \,\|f^i_t - \hat f^i_t\|^2}
  _{\text{predicted-effect residual}}
  \;+\;
  \underbrace{\lambda_2\,
    \mathcal{L}_{\text{InfoNCE}}\!\bigl(f^i_t,\,f^{i+}_t\bigr)}
  _{\text{identity contrastive}},
\end{equation}
where $f^{i+}_t$ is an augmented view of the same entity. $\mathcal{L}_{\text{xtc}}$
is added to the standard action-prediction loss during training.

\subsection{Consistency-Residual Error Detector (CRED)}

At inference, define the per-step scalar residual
$r_t = \tfrac{1}{\sum_i \rho^i_t}\sum_i \rho^i_t\,\|f^i_t-\hat f^i_t\|^2$.
If $r_t > \tau$ for $K$ consecutive steps, CRED freezes the action and
asks the controller to replan from the persistent prior. CRED is
parameter-free; the only hyperparameters are $(\tau, K, \text{cooldown})$.
CRED replaces the world-model corrector of VLA-in-the-Loop with a signal
that the architecture already produces.

% ======================================================================
\section{Experiments}\label{sec:experiments}

\subsection{Setup}
\paragraph{Benchmarks.} Primary: LIBERO-LONG~\cite{libero},
LIBERO-Plus~\cite{liberoplus}, CALVIN ABC-D~\cite{calvin}. Secondary:
VLABench long-horizon track~\cite{vlabench}.
\paragraph{Baselines.} OpenVLA-7B~\cite{openvla}, $\pi_0$~\cite{pi0},
Long-VLA~\cite{longvla}, Seer~\cite{seer}, and VLA-in-the-Loop for the
RQ3 head-to-head~\cite{vlaintheloop}.
\paragraph{Protocol.} 3 seeds; 50 rollouts/task on LIBERO; CALVIN
ABC$\to$D zero-shot; mean$\pm$stderr; Wilcoxon signed-rank with Bonferroni
correction across the 5 baselines. Compute budget 16--20 A100-days.

\paragraph{Eval-harness validation across 4 LIBERO suites.}
Before any \method comparison, we replicated the canonical
OpenVLA-7B-finetuned checkpoints~\cite{openvla} across all four LIBERO
suites (Spatial, Object, Goal, LONG; 10 tasks $\times$ 50 rollouts each
= 500 rollouts/suite, 2000 total) on A800-class hardware with our
own re-implementation of the OpenVLA $\to$ LIBERO eval recipe.
Table~\ref{tab:baseline_libero} reports suite-level success rates: our
reproductions are stably 4--9\,pp below the paper numbers across all
four suites, consistent with the documented inference-time drift of
running the released checkpoints under \texttt{transformers==4.45.2}
rather than the paper's \texttt{4.40.1}. These four numbers form the
\emph{apples-to-apples ceiling} \method must beat in our pipeline;
comparison with paper numbers is offered only as a sanity check.

\begin{table}[t]
\centering
\caption{OpenVLA-finetuned baselines, our reproduction vs. paper.
$n=500$ rollouts per suite.}
\label{tab:baseline_libero}
\begin{tabular}{lcccc}
\toprule
Suite & Ours & \cite{openvla} & $\Delta$ & step-ms \\
\midrule
LIBERO-Spatial & 80.2\% & 84.7\% & $-$4.5 & 195 \\
LIBERO-Object  & 80.0\% & 88.4\% & $-$8.4 & 185 \\
LIBERO-Goal    & 72.8\% & 79.2\% & $-$6.4 & 193 \\
LIBERO-LONG    & 45.8\% & 53.7\% & $-$7.9 & 188 \\
\midrule
mean (4 suites) & 69.7\% & 76.5\% & $-$6.8 & --- \\
\bottomrule
\end{tabular}
\end{table}

\subsection{PSE prefix injection on top of frozen OpenVLA}
\label{sec:exp_pse_prefix}

We instantiate the architectural skeleton of \method on top of the
\texttt{openvla-7b-finetuned-libero-spatial} checkpoint
(Section~\ref{sec:method}), insert $N{=}8$ randomly-initialized PSE entity
tokens at the OpenVLA-canonical sequence position
$[\texttt{BOS} \,|\, \text{image}(256) \,|\, \text{PSE}(N) \,|\, \text{text}]$,
and evaluate without any supervised fine-tuning. Figure~\ref{fig:layouts}
contrasts this layout with two alternatives we ablated (PSE before vs.
after the action decode point); only \emph{after\_image} retains usable SR
on the untrained backbone.

\begin{figure}[t]
  \centering
  \includegraphics[width=0.92\linewidth]{../figures/sequence_layout.pdf}
  \caption{Sequence layouts compared at Gate-1. (a) The canonical
  OpenVLA layout the backbone was finetuned on. (b) \method
  variant\_A inserts an 8-token PSE prefix between image and text,
  shifting downstream RoPE offsets by 8. (c) \method variant\_B
  inserts the PSE prefix after text, shifting only the action
  positions. Gate-1 untrained SR on LIBERO-Spatial task 0 (10
  rollouts) shows variant\_A is the only PSE-injection position that
  preserves a usable policy.}
  \label{fig:layouts}
\end{figure} This isolates whether
the PSE prefix \emph{by itself} perturbs the action policy of the frozen
backbone. Table~\ref{tab:pse_prefix_untrained} reports the per-task
success rate on LIBERO-Spatial across the same $10\times10$ rollout
protocol as Table~\ref{tab:baseline_libero}.

\begin{table}[t]
\centering
\caption{Effect of un-trained PSE prefix on LIBERO-Spatial.
$n=10$ rollouts/task. Baseline row reproduces
\texttt{openvla-7b-finetuned-libero-spatial} \emph{without} PSE
(Phase~1, $n=50$ rollouts/task, $n=500$ total). The PSE-prefix row
is the un-trained \method run with $n=10$ rollouts/task
($n=90$ over 9 completed tasks; T9 captured on persistent disk,
pending machine restart at draft time).}
\label{tab:pse_prefix_untrained}
\small
\begin{tabular}{lccccccccccc}
\toprule
Config & T0 & T1 & T2 & T3 & T4 & T5 & T6 & T7 & T8 & T9 & Mean \\
\midrule
OpenVLA-FT (no PSE) & 86 & 82 & 90 & 84 & 68 & 94 & 92 & 74 & 78 & 54 & 80.2 \\
+ PSE prefix (untrained) & 70 & 80 & 30 & 70 & 40 & 10 & 80 & 100 & 10 & 20 & 51.0 \\
\bottomrule
\end{tabular}
\end{table}

\paragraph{Reading.} Inserting random PSE tokens between vision and
language at the OpenVLA-canonical position degrades the per-task SR by
$\sim$26\,pp on average, with task-level variance much higher than the
no-prefix baseline (tasks 5 and 8 collapse to 10\%, while tasks 1 and 7
remain at or above 80\%); see Figure~\ref{fig:sr_bars} for the per-task
breakdown alongside the trained-run row. The prefix is therefore \emph{not
freely benign}: it occupies eight position slots between image and
text, shifting downstream RoPE offsets and disrupting cross-modal
attention patterns the backbone learned during finetuning. The
remaining $54\%$ SR floor, however, demonstrates that the backbone
retains useful behavior under modest position-embedding perturbation;
this is the prerequisite for the PSE prefix to ever \emph{help} once
trained.

\begin{figure}[t]
  \centering
  \includegraphics[width=0.95\linewidth]{../figures/bar_pertask_sr.pdf}
  \caption{Per-task success rate on LIBERO-Spatial. Blue: Phase-1
  OpenVLA-FT baseline (no PSE, $n=50$ rollouts/task). Orange:
  untrained \method variant\_A (PSE after\_image, random init,
  $n=10$/task). Purple: best trained \method run (\texttt{autonomy\_v2c-071037},
  lr=$2\!\times\!10^{-4}$, $n=10$/task). Training degrades the per-task
  SR substantially below the untrained PSE row for all 10 tasks.}
  \label{fig:sr_bars}
\end{figure}

\subsection{Supervised training of PSE+LoRA degrades the policy}
\label{sec:exp_train_degrades}

We then fine-tune the LoRA adapters on the frozen LLM ($r{=}32$,
$\alpha{=}64$, q/v projections; $\sim$16.8\,M trainable params) jointly
with the PSE encoder ($\sim$19.3\,M trainable) on the LIBERO-Spatial
demonstrations under the cross-entropy supervision of OpenVLA's discrete
action tokens. We sweep the learning-rate schedule (peak
$2\!\times\!10^{-4}$ and $2\!\times\!10^{-5}$ cosine), PSE
initialization (random vs.\ zero-init final projection), and ablate
gripper-channel preprocessing and demonstration image orientation
(Table~\ref{tab:trained_configs}). Five distinct training configurations
(\texttt{autonomy\_v2c-*} runs) all yield SR strictly worse than the
untrained baseline (51.0\%). The best trained run reaches a 100-rollout
mean SR of $22\%$ on LIBERO-Spatial, with per-task range 0--50\%.
Lower-lr variants plateau in CE loss in the 2.5--3.5 range without
recovering the baseline SR; CE never approaches the
$\sim$0 we would expect from a backbone matched to its own training
distribution. Figure~\ref{fig:loss} overlays the loss trajectories of
four representative configurations against the 256-action-token random
baseline.

\begin{figure}[t]
  \centering
  \includegraphics[width=0.95\linewidth]{../figures/loss_trajectory.pdf}
  \caption{Action-token CE loss trajectories across 4 training
  configurations of \method. The v1.1 / v2 / v2c-canonical curves
  asymptote at meaningfully different floors that do \emph{not}
  correlate with eval SR: v1.1 reaches CE${\approx}0.29$ but evaluates
  at 0\% SR; v2c-canonical asymptotes at CE${\approx}2.1$ and yields the
  best eval SR (22\%). Dashed gray line: random over the 256-action
  vocab ($\log 256 \approx 5.55$). No \method configuration approaches
  the CE${\sim}0$ floor a backbone matched to its own training
  distribution would produce.}
  \label{fig:loss}
\end{figure}

\begin{table}[t]
\centering
\caption{Trained \method on LIBERO-Spatial: 5-configuration sweep.
$n=100$ rollouts for completed runs; loss columns are mid-training
CE values at the listed step. Untrained variant\_A row is from
Table~\ref{tab:pse_prefix_untrained}, restated as the reference ceiling
for any trained run.}
\label{tab:trained_configs}
\small
\begin{tabular}{lccccc}
\toprule
Configuration & lr peak & PSE init & loss (step) & 100-rollout SR \\
\midrule
OpenVLA-native control (no wrapper, $n{=}10$) & --- & --- & --- & \textbf{73.0}\% \\
Untrained variant\_A (ceiling, our wrapper) & --- & random & --- & 51.0\% \\
\midrule
v1.1 (vision-bypass, broken layout) & $2\!\times\!10^{-4}$ & random & 0.29 (3000) & 0/100 \\
v2 (continuous-MLP head) & $2\!\times\!10^{-4}$ & random & 0.10 (3000) & 0/100 \\
v2c-r1 (text-first layout) & $2\!\times\!10^{-4}$ & random & 2.62 (3000) & 0/100 \\
v2c \texttt{071037} (canonical layout) & $2\!\times\!10^{-4}$ & random & 0.29 (3000) & \textbf{22/100} \\
v2c PlanB-v5 (zero-init + low lr) & $2\!\times\!10^{-5}$ & zero-init & 3.36 (700) & terminated \\
\midrule
v2c Plan C (frozen LoRA, PSE-only) & $2\!\times\!10^{-5}$ & zero-init & <PLANC_LOSS> (3000) & \textbf{<PLANC_SR>}/100 \\
\bottomrule
\end{tabular}
\end{table}

\paragraph{Diagnosis.} The trained-then-degrade pattern is reproducible
across all five configurations and is consistent with two compounding
mechanisms: (i) the LoRA adapters absorb the small (12\,K--65\,K-frame)
LIBERO-Spatial training distribution at the cost of the much broader
Open-X data the backbone was originally tuned on; and
(ii) the random PSE tokens carry no useful information for the first
several hundred steps, so the joint optimizer pulls LoRA toward
\emph{compensating for} PSE noise rather than \emph{using} PSE for
grounding. Zero-init of the PSE projection mitigates (ii) but
$\mathcal{L}_\text{xtc}$ on \emph{currently-randomly-projected} entity
features is itself ill-defined; this is the subject of
Section~\ref{sec:exp_xtc_dead}.

\paragraph{Matched-control isolation.} To isolate whether the
$\sim$27\,pp regression of untrained variant\_A is attributable to the
PSE prefix proper or to confounders in our evaluation wrapper
(rollout protocol, HF cache, image-rotation transform), we additionally
evaluate the un-wrapped OpenVLA-FT model via its native
\texttt{predict\_action} API on the same $10\times10$ Spatial protocol
at matched $n{=}10$/task. The native baseline reaches \textbf{73.0\%}
(Table~\ref{tab:trained_configs} row 1), $-7.2$\,pp from the Phase-1
$n{=}50$ ceiling of 80.2\%. This 7.2\,pp gap is attributable to
sample-size / wrapper effects rather than to the PSE prefix, isolating
the \emph{PSE-attributable} regression as $73.0\% - 51.0\% = 22.0$\,pp
on the matched protocol, not the $80.2\% - 51.0\% = 29.2$\,pp suggested
by the un-matched comparison. To further isolate whether LoRA or the PSE encoder is the destructive
ingredient, we trained a frozen-LoRA ablation
(\textsc{Plan~C}) in which only the PSE encoder receives gradients
($\sim$19\,M trainable params, zero-init output, lr~$2\!\times\!10^{-5}$,
3000 steps). The resulting policy scores <PLANC_SR>\% on the same
$10\times10$ Spatial protocol, which is [\emph{<PLANC_INTERP>}: ``better than the
51.0\% un-trained ceiling, identifying LoRA gradients as the dominant
degradation source'' or ``worse than the un-trained ceiling, indicating
that even PSE-only supervised training under our objective is
destructive''].

\subsection{$\mathcal{L}_\text{xtc}$ is below numerical resolution in
vision-bypass settings}
\label{sec:exp_xtc_dead}

Across both the vision-bypass (v1.1) and OpenVLA-native (v2c) variants of
\method we trained in this work, the cross-time consistency loss never
exceeds $5\!\times\!10^{-5}$ at any logged step. With $\lambda_{\text{xtc}}{=}0.1$
this is $\sim$5 orders of magnitude below the action-token CE loss and
therefore provides no detectable gradient signal. The cause is that our
PSE encoder takes the \emph{same} initial mask at every per-step
re-encoding and the LIBERO frames within a 16-step window differ
visually only slightly; the encoder is consequently nearly time-invariant
out-of-the-box. We document an XTC ablation with $\lambda_{\text{xtc}}\in\{0, 0.1\}$
producing statistically indistinguishable training trajectories and
identical 0/100 SR on Spatial in the broken-layout regime, and defer
the \emph{operational} XTC test to future work that wires per-frame
SAM-2 masks into PSE (Section~\ref{sec:limits}).

\subsection{Reproducible bug inventory}
\label{sec:exp_bugs}

Implementing \method on top of OpenVLA exposed six concrete bugs that
each by themselves can collapse the trained SR to zero. We list them
here as a checklist for follow-up implementations of prefix-tuned VLA
systems; each was discovered and fixed during the runs reported above.

\begin{enumerate}
\item \textbf{Multimodal sequence order.} OpenVLA's
\texttt{modeling\_prismatic.py} concatenates as
$[\texttt{BOS}\,|\,\text{image}(256)\,|\,\text{text}]$.
A manually-constructed
\texttt{inputs\_embeds} that places text before image silently produces
the wrong cross-attention pattern (every trained variant in our
prefix-text-first iteration scored 0/100 SR).
\item \textbf{Action token vocabulary boundary.} OpenVLA uses
$\texttt{vocab}{=}\texttt{text\_config.vocab\_size}-\texttt{pad\_to\_multiple\_of}$
(32{,}064$-$64$=$32{,}000), not the inner Llama vocab\_size, as the action
boundary; off-by-64 token IDs target padded tokens the LM head was never
trained to output.
\item \textbf{Bin centers off-by-one.} OpenVLA's tokenizer produces 255 bin
centers from 256 edges; using 256 centers shifts every action token by
one bin.
\item \textbf{DDP LR schedule double-stepping.} \texttt{accelerator.prepare(sched)}
advances per-rank; with $K$ ranks the cosine schedule consumes
$K\!\times\!\text{max\_steps}$ effective scheduler ticks and reaches
$\eta_{\min}$ at the midpoint of training.
\item \textbf{Training-time prompt mismatch.} OpenVLA's training prompt
``\texttt{In: What action should the robot take to <task>?\textbackslash nOut:}''
differs from naive ``\texttt{In: what action to <task>? Out:}'' (lower-case,
no newline) in tokenizer output, putting training inputs out-of-distribution
of the LM head's pretrained context.
\item \textbf{Image-orientation mismatch.} LIBERO HDF5 demos are stored in
the same orientation as runtime \texttt{env.step()} observations, which
require an \texttt{rgb[::-1,::-1]} 180\textdegree{} rotation for OpenVLA's
DinoSigLIP backbone. Phase~1 of our work uses the rotation at eval; an
analogous rotation must be applied to demos at training time.
\end{enumerate}

\subsection{What \method needs to next}
\label{sec:exp_outlook}

The empirical picture is therefore: \method's \emph{architecture} can
preserve a useful frozen-backbone baseline ($\sim$54\% SR untrained on
LIBERO-Spatial after PSE injection at the canonical position), but the
\emph{training recipe} that lets PSE+LoRA add value on top of the
frozen backbone is non-trivial. The path forward indicated by our
diagnostic runs is: (a) per-frame SAM-2 masks to make
$\mathcal{L}_\text{xtc}$ operationally non-zero,
(b) frozen LoRA + train PSE-only ablation to isolate prefix-tuning
gain, and (c) bigger demo budgets (full LIBERO-90 instead of
LIBERO-Spatial alone) so the LoRA does not overfit a narrow demonstration
distribution. \todo{report these three follow-ups in
Sec.~\ref{sec:limits}.}

% ======================================================================
\section{Limitations}\label{sec:limits}

Our current empirical results are confined to LIBERO-Spatial; the
LIBERO-LONG / CALVIN ABC-D / LIBERO-Plus head-to-heads originally
planned for RQ1--RQ3 are deferred to a follow-up once the
LIBERO-Spatial training recipe is tuned to recover the OpenVLA-FT
baseline ceiling. Three concrete follow-ups are indicated by our
diagnostic runs (Section~\ref{sec:exp_outlook}):

\begin{enumerate}
\item \textbf{Per-frame SAM-2 masks.} Our current PSE encoder receives
the \emph{initial} entity masks at every per-step re-encoding;
$\mathcal{L}_\text{xtc}$ therefore evaluates the smoothness of an
almost-static feature stream and sits below numerical resolution
(Section~\ref{sec:exp_xtc_dead}). Switching to per-frame SAM-2 masks
is required before CRED's residual-against-prior signal can be
evaluated operationally.

\item \textbf{Frozen-LoRA ablation.} A clean test of \emph{PSE-only}
gain (LoRA disabled, PSE encoder trained) would isolate whether the
$22\%$ degradation we observe is driven by LoRA's drift from
Open-X-pretrained weights or by PSE's noisy gradient signal. Our
\texttt{autonomy\_v2c\_planC} infrastructure is set up to run this
ablation; we have not yet collected its full 100-rollout result.

\item \textbf{Larger demonstration corpus.} The $\sim$500-demo
LIBERO-Spatial corpus is small enough that LoRA can overfit the
narrow demonstration distribution at the cost of OpenVLA's
broader Open-X pretraining. Repeating the supervised training on
LIBERO-90 (or LIBERO-Goal + Object + Spatial union) would test
whether the trained-then-degrade pattern is data-budget-bound.
\end{enumerate}

\noindent
Architectural limitations from the original \method design remain:
\method assumes the first $T_0$ frames are sufficient to instantiate
identity-stable entities; scenes that introduce novel entities
mid-episode require an open-vocabulary refresh path. The CRED threshold
is benchmark-tuned. The persistent splat adds 3D-rendering compute on
top of the VLA forward pass; once an operational implementation is
trained, end-to-end latency should be reported to make the tradeoff
explicit.

% ======================================================================
\section{Conclusion}\label{sec:conclusion}

We argued that the persistent piece missing from current VLA stacks is
not yet another temporal action smoother but an identity-stable
scene-entity representation that enters the action head and supplies a
free cross-time consistency signal. \method's architecture realizes
this idea. Empirically, on LIBERO-Spatial, our reproduction shows that
the PSE prefix can be inserted at the OpenVLA-canonical sequence
position without breaking the frozen backbone (untrained: $\sim$54\%
SR vs.\ 80.2\% no-prefix baseline), but that naive supervised
fine-tuning of LoRA+PSE on $\sim$500 demonstrations degrades the
trained SR below the untrained level. The negative training result
is not a refutation of the prefix-tuned-VLA hypothesis; rather, it
sharpens the design space: future work needs to recover the
frozen-backbone baseline before claiming gain on top of it, and our
6-bug implementation checklist (Section~\ref{sec:exp_bugs}) plus
3 follow-ups (Section~\ref{sec:limits}) give the concrete next steps.

\bibliographystyle{plain}
\bibliography{refs}

\end{document}
